% Source: source_md/intelligence_part2_appendices_ja.md
\section{領域横断ホモロジー表}


本補遺は、VED 理論ツリー全体にわたる領域横断ホモロジーを表として整理する。本論文の主要概念が、他の層・他の論文でいかなる構造的相同物を持つかを示す。

\bigskip


\subsection{停止・閉包・アトラクターのホモロジー}


\begin{longtable}{@{}p{0.18\linewidth}p{0.18\linewidth}p{0.18\linewidth}p{0.18\linewidth}p{0.18\linewidth}@{}}
\toprule
知性 Part II(本論文) & Vol.1〜4(物理層) & Bio-IFGT(生命層) & 意識編 & 概念地形編 \\
\midrule
\endfirsthead
\toprule
知性 Part II(本論文) & Vol.1〜4(物理層) & Bio-IFGT(生命層) & 意識編 & 概念地形編 \\
\midrule
\endhead
内部不動点不在($L_F$ の非閉包) & $C_{max}$ 不到達性 & 準閉包の維持限界 & 意識の状態遷移が完結しない & 概念の完全閉包は到達不能 \\
創発停止層 $L_E$(アトラクター) & 粒子 = $E(A)$ の局所極小(Vol.2) & 生命アトラクター $\Omega_{\text{life}}$ & 意識の安定状態 & 概念の準安定点 \\
第二アトラクター($L_E$) & ゲージ群の安定構造 & 第二アトラクター $\Omega_{\text{intel}}$ & — & — \\
差分の地平線(知性層) & 差分の地平線 $C_{max}$(Vol.4) & 生命の非平衡限界 & 意識の差分地平 & 概念地形の境界 \\
分岐爆発(推論の) & 場の発散 & 生命の分岐構造 & 意識状態の分岐 & 概念分岐の指数増殖 \\
\bottomrule
\end{longtable}


\bigskip


\subsection{創発構造のホモロジー($\mathcal{U} = \Phi(\mathcal{L}, H_{ij})$)}


\begin{longtable}{@{}p{0.18\linewidth}p{0.18\linewidth}p{0.18\linewidth}p{0.18\linewidth}p{0.18\linewidth}@{}}
\toprule
創発例 & 下位層 $\mathcal{L}$ & 外部相互拘束 $H_{ij}$ & 上位層 $\mathcal{U}$ & 参照 \\
\midrule
\endfirsthead
\toprule
創発例 & 下位層 $\mathcal{L}$ & 外部相互拘束 $H_{ij}$ & 上位層 $\mathcal{U}$ & 参照 \\
\midrule
\endhead
標準模型構造 & 因果ログ $C_{ij}$・内部ベクトル $\vec{e}_i$ & 閉包条件、閉包エネルギー $E(A)$ & 三角形閉包・粒子(局所極小) & Vol.2 \\
生命 & 情報流 $J$・状態発展 $\dot{x}=F(x;\theta)$ & DNA 由来パラメータ $\theta$・非平衡交換 & 生命アトラクター $\Omega_{\text{life}}$ & Bio-IFGT \\
創発停止層 & ログ層 $L_R$ の再帰運動 & 観測者-環境相互拘束 $H_{ij}$ & 創発停止層 $L_E$(停止アトラクター) & 本論文 \S{}5 \\
\bottomrule
\end{longtable}


すべての例が定義 1.1 の形式 $\mathcal{U} = \Phi(\mathcal{L}, H_{ij})$ を満たす。VED 理論ツリーはアトラクター構造の創発を共通言語として使う。

\bigskip


\subsection{projection(投影)のホモロジー}


\begin{longtable}{@{}p{0.30\linewidth}p{0.30\linewidth}p{0.30\linewidth}@{}}
\toprule
層 & projection の内容 & 参照 \\
\midrule
\endfirsthead
\toprule
層 & projection の内容 & 参照 \\
\midrule
\endhead
意識編(IFGT 第四プリミティブ) & 状態・表現が不均等な影響力(質量)を獲得する操作。$\Phi$ と $J_I$ を結合。主体を前提しない & 意識編 \texttt{formal\_mapping} \\
知性 Part II(\S{}5.5〜\S{}5.6) & (B) 組み替えによる停止演算子設置領域の投影。外部相互拘束の構造が状態空間上に写される & 本論文 \S{}5.5 \\
心理学的 projection & 内的状態を外部対象に帰属する。主体・無意識・防衛機制を前提 & (比較のみ・本論文の主張ではない) \\
\bottomrule
\end{longtable}


作用的近接: 意識編と知性 Part II の投影は「内部で閉じないものを外部に投射して処理する」という作用形式を共有しつつ、主体を前提しない点で心理学的 projection と区別される。

\bigskip


\subsection{$\theta$(拘束パラメータ)のホモロジー}


\begin{longtable}{@{}p{0.30\linewidth}p{0.30\linewidth}p{0.30\linewidth}@{}}
\toprule
層 & $\theta$ の実体 & 可変時間スケール \\
\midrule
\endfirsthead
\toprule
層 & $\theta$ の実体 & 可変時間スケール \\
\midrule
\endhead
Bio-IFGT & DNA・反応傾向・閾値 & 世代スケール \\
本論文(知性一般) & 外部条件・環境設定 & 環境変化スケール \\
本論文(サピエンス) & 推論操作の拘束・履歴的閾値 & 個体内・対話内スケール \\
\bottomrule
\end{longtable}


同一役割(運動可能領域を画定し生成に構造を与える)を果たしつつ、可変の時間スケールが層によって異なる。

\bigskip


\subsection{F(状態発展演算子)のホモロジー}


\begin{longtable}{@{}p{0.30\linewidth}p{0.30\linewidth}p{0.30\linewidth}@{}}
\toprule
層 & $F$ の意味 & 参照 \\
\midrule
\endfirsthead
\toprule
層 & $F$ の意味 & 参照 \\
\midrule
\endhead
Bio-IFGT & 生命系の状態発展演算子 $\dot{x} = F(x;\theta)$ & Bio-IFGT \texttt{core\_equations} \\
本論文 知性層 & 推論演算子 $x_{n+1} = F(x_n)$。$F(x;\theta)$ の知性層特殊形 & 本論文 \S{}2.1 \\
\bottomrule
\end{longtable}


同じ記号 $F$ が異なる層で同型の役割を果たす。これは記号の衝突ではなく、理論的連続性の証拠として理解される。

\bigskip


\subsection{停止構造の歴史的形態ホモロジー}


本論文が分析した停止構造の形態と、VED の他概念との対応:

\begin{longtable}{@{}p{0.18\linewidth}p{0.18\linewidth}p{0.18\linewidth}p{0.18\linewidth}p{0.18\linewidth}@{}}
\toprule
停止形態 & $\Phi_{\text{net}}$ 構造 & $L_E$ の維持様態 & $\mu_\theta$ の関与 & 時間スケール \\
\midrule
\endfirsthead
\toprule
停止形態 & $\Phi_{\text{net}}$ 構造 & $L_E$ の維持様態 & $\mu_\theta$ の関与 & 時間スケール \\
\midrule
\endhead
粘菌的停止 & 環境(スライム跡+物理的応答) & 環境相互干渉で自律創発 & 固定(外部受動) & 個体の行動時間 \\
アニミズム & 環境的・分散的 & 環境直接相互干渉 & 最初の社会的具現化 & 個体・小集団時間 \\
神 & 中心化された単一参照点 & 単一懸架軌跡 & 中心化による安定化 & 文明スケール \\
宗教 & 中心化+分散の混合相 & 世代横断的維持 & 世代横断的安定化 & 世代・文明スケール \\
社会的BC & 分散(個体間 $L_E$ ネット) & 分散ノード間同期 & 分散的可動性 & 制度・歴史スケール \\
Transformer & 外部注入(タスク仕様) & 外部注入 & 外部設定 & タスク時間 \\
サピエンス & 個体内+社会共有 & 個体内創発+社会同期 & 内発的(短期) & 対話〜生涯 \\
\bottomrule
\end{longtable}


\bigskip


\subsection{Part I との構造対応表}


\begin{longtable}{@{}p{0.30\linewidth}p{0.30\linewidth}p{0.30\linewidth}@{}}
\toprule
本論文(Part II) & 知性編 Part I & 対応の性格 \\
\midrule
\endfirsthead
\toprule
本論文(Part II) & 知性編 Part I & 対応の性格 \\
\midrule
\endhead
非閉包定理 \S{}3 & \S{}8 個体知性の非閉包 & Part II が Part I を定理として証明 \\
(A)/(B) 棲み分け & — & Part II の新規追加。Part I \S{}8 を精密化 \\
差分の地平線(知性層) & \S{}9 知性の差分の地平線 & 同一構造の別記述 \\
ゲティア問題と $\mu_\theta$ の同値 & \S{}6 ゲティアの構造的診断 & Part II が $\mu_\theta$ との同値を追加 \\
三層プロファイル比較 & \S{}5 Physarum・Transformer & Part II が三層で再位置づけ \\
分岐爆発 = 内部不動点不在の動力学 & \S{}7 動力学的不安定性 & 同一現象の二側面として統合 \\
$L_E$ = 第二アトラクターの停止構造 & — & Bio-IFGT の第二アトラクターを Part II が展開 \\
\bottomrule
\end{longtable}


\bigskip
